\SourceLeaf{162}
Facile vero perspicitur, per eandem transformationem, per quam $(\frac{A}{m},\frac{B}{m},\frac{C}{m})$ transeat in $(\pm1,0,\pm1)$, formam $(A,B,C)$ transire in $(\pm m,0,\pm m)$, ancipitem. Quare forma $(A,B,C)$. ancipiti aequivalens, cuivis formae, cui aequivalet, tum proprie tum improprie aequivalebit.

3) Si $\frac{4D}{mm}=3$, sive $4D=3mm$. Tum $m$ erit par omnesque solutiones aequationis $tt+Duu=mm$ erunt sex,
\[
t,u;=m,0;\;-m,0;\;\tfrac12m,1;\;-\tfrac12m,-1;\;\tfrac12m,-1;\;-\tfrac12m,1
\]
Si itaque duae transformationes dissimiles formae $F$ in $f$ habentur,
\[
\begin{aligned}
x&=\alpha x'+\beta y',& y&=\gamma x'+\delta y'\\
x&=\alpha'x'+\beta'y',& y&=\gamma'x'+\delta'y'
\end{aligned}
\]
habebuntur duodecim transformationes, scilicet sex priori similes
\[
x=\pm\alpha x'\pm\beta y',\quad y=\pm\gamma x'\pm\delta y'
\]
\[
\begin{aligned}
x&=\pm(\tfrac12\alpha-\tfrac{\alpha B+\gamma C}{m})x'\pm(\tfrac12\beta-\tfrac{\beta B+\delta C}{m})y'\\
y&=\pm(\tfrac12\gamma+\tfrac{\alpha A+\gamma B}{m})x'\pm(\tfrac12\delta+\tfrac{\beta A+\delta B}{m})y'
\end{aligned}
\]
\[
\begin{aligned}
x&=\pm(\tfrac12\alpha+\tfrac{\alpha B+\gamma C}{m})x'\pm(\tfrac12\beta+\tfrac{\beta B+\delta C}{m})y'\\
y&=\pm(\tfrac12\gamma-\tfrac{\alpha A+\gamma B}{m})x'\pm(\tfrac12\delta-\tfrac{\beta A+\delta B}{m})y'
\end{aligned}
\]
et sex posteriori similes, quae ex his nascuntur ponendo pro $\alpha$, $\beta$, $\gamma$, $\delta$ hos $\alpha'$, $\beta'$, $\gamma'$, $\delta'$.

Quod vero in hoc casu semper $F,f$ utroque modo aequivalent, ita demonstramus. Formae $(\frac{2A}{m},\frac{2B}{m},\frac{2C}{m})$ determinans erit $=-\frac{4D}{mm}=-3$. adeoque haec forma (art. 176) aut formae $(\pm1,0,\pm3)$ aut huic $(\pm2,\pm1,\pm2)$ aequivalens. Unde facile perspicitur, formam $(A,B,C)$ aut formae $(\pm\tfrac12m,0,\pm\tfrac32m)$ aut huic $(\pm m,\tfrac12m,\pm m)$\HistoricalNote{*)}{Demonstrari potest, formam $(A,B,C)$ necessario posteriori aequivalere: sed hoc hic non necessarium.} quae ambae sunt ancipites, aequivalere adeoque, cuivis aequivalenti, utroque modo.

4) Si supponitur $\frac{4D}{mm}=2$, fit $(\frac{2B}{m})^2=4\frac{AC}{mm}-2$, adeoque $\equiv2\ (\mathrm{mod.}\ 4)$.
\SourceLeaf{163}
Sed quum nullum quadratum esse possit $\equiv2\ (\mathrm{mod.}\ 4)$, hic casus locum habere nequit.

5) Supponendo $\frac{4D}{mm}=1$, fit $(\frac{2B}{m})^2=4\frac{AC}{mm}-1\equiv-1\ (\mathrm{mod.}\ 4)$. Quod quum impossibile sit, etiam hic casus nequit locum habere.

Ceterum quum $D$ neque $=0$, neque negativus sit, alii casus praeter enumeratos dari non possunt.

\Needspace{8\baselineskip}
\Article{180}
\textsc{Problema.} \textit{Invenire omnes repraesentationes numeri dati $M$ per formam $axx+2bxy+cyy\ldots F$, determinantis negativi $-D$, in quibus $x,y$ valores inter se primos nanciscantur.}

\textit{Sol.} Ex art. 154 patet, $M$ eo quo requiritur modo repraesentari non posse, nisi $-D$ sit resid. quadr. ipsius $M$. Investigentur itaque primo omnes valores diversi (\textit{i. e.} incongrui) expr. $\surd-D\ (\mathrm{mod.}\ M)$, qui sint $N,-N,N',-N',N'',-N''$ etc.; quo simplicior evadat calculus, omnes $N,N'$ etc. ita determinari possunt, ut non sint $>\tfrac12M$. Iam quoniam quaevis repraesentatio ad aliquem horum valorum pertinere debet, singuli seorsim considerentur.

Si formae $F,(M,N,\frac{D+NN}{M})$ non sunt proprie aequivalentes, nulla repraesentatio ipsius $M$ ad valorem $N$ pertinens dari potest (art. 168). Si vero sunt, investigetur transformatio propria formae $F$ in
\[
Mx'x'+2Nx'y'+\frac{D+NN}{M}y'y'
\]
quae sit
\[
x=\alpha x'+\beta y',\quad y=\gamma x'+\delta y'
\]
eritque $x=\alpha,y=\gamma$ repraesentatio numeri $M$ per $F$ ad $N$ pertinens. Sit div. comm. max. numerorum $A,2B,C,=m$ distinguanturque tres casus (art. praec.):

1) Si $\frac{4D}{mm}>4$, aliae repraesentationes ad $N$ pertinentes quam hae \textit{duae} $x=\alpha,y=\gamma$; $x=-\alpha,y=-\gamma$ non dabuntur (artt. 169, 179).

2) Si $\frac{4D}{mm}=4$, habebuntur \textit{quatuor} repraesentationes
\[
x=\pm\alpha,\quad y=\pm\gamma;\quad x=\mp\frac{\alpha B+\gamma C}{m},\quad y=\pm\frac{\alpha A+\gamma B}{m}
\]
3) Si $\frac{4D}{mm}=3$, habebuntur \textit{sex} repraesentationes
\SourceLeaf{164}
\[
\begin{aligned}
x&=\pm\alpha,&y&=\pm\gamma;\\
x&=\pm(\tfrac12\alpha-\tfrac{\alpha B+\gamma C}{m}),&y&=\pm(\tfrac12\gamma+\tfrac{\alpha A+\gamma B}{m})\\
x&=\pm(\tfrac12\alpha+\tfrac{\alpha B+\gamma C}{m}),&y&=\pm(\tfrac12\gamma-\tfrac{\alpha A+\gamma B}{m})
\end{aligned}
\]
Eodem modo quaerendae sunt repraesentationes ad valores $-N.N',-N'$ etc. pertinentes.

\Needspace{8\baselineskip}
\Article{181}
Investigatio repraesentationum numeri $M$ per formam $F$, in quibus $x,y$ valores inter se non primos habent, ad casum iam consideratum facile reduci potest. Fiat talis repraesentatio ponendo $x=\mu e,y=\mu f$, ita ut $\mu$ sit div. comm. max. ipsorum $\mu e,\mu f$, sive $e,f$ inter se primi. Tum erit $M=\mu\mu(Aee+2Bef+Cff)$ adeoque per $\mu\mu$ divisibilis; substitutio vero $x=e,y=f$ erit repraesentatio numeri $\frac{M}{\mu\mu}$ per formam $F$, in qua $x,y$ valores inter se primos habent. Si itaque $M$ per nullum quadratum (praeter 1) divisibilis est, \textit{e. g.} si est numerus primus: tales repraesentationes ipsius $M$ non dabuntur. Si vero $M$ divisores quadraticos implicat, sint hi $\mu\mu,\nu\nu,\pi\pi$, etc. Quaerantur primo omnes repraesentationes numeri $\frac{M}{\mu\mu}$ per formam $(A,B,C)$, in quibus $x,y$ valores inter se primos habent, qui valores si per $\mu$ multiplicantur, praebebunt omnes repraesentationes ipsius $M$, in quibus div. comm. max. numerorum $x,y$ est $\mu$. Simili modo omnes repraesentationes ipsius $\frac{M}{\nu\nu}$, in quibus valores ipsorum $x,y$ inter se primi sunt, praebebunt omnes repraesentationes ipsius $M$, in quibus div. comm. max. valorum ipsorum $x,y$ est $\nu$ etc.

Palam igitur est, per praecepta praecedentia omnes repraesentationes numeri dati per formam datam determinantis negativi inveniri posse.

\Needspace{12\baselineskip}
\begin{center}
\textit{Applicationes speciales ad discerptionem numerorum in quadrata duo, in quadratum simplex et duplex, in simplex et triplex.}
\end{center}

\Needspace{8\baselineskip}
\Article{182}
Descendimus ad quosdam casus particulares, tum propter insignem ipsorum elegantiam tum propter assiduam operam ab ill. Eulero ipsis impensam, unde classicam quasi dignitatem sunt nacti.

I. Per formam $xx+yy$ ita repraesentari ut $x$ ad $y$ sit primus (sive in
\SourceLeaf{165}
duo quadrata inter se prima discerpi) nullus numerus potest nisi cuius residuum quadraticum est $-1$, tales vero numeri, positive accepti, omnes poterunt. Sit $M$ talis numerus, omnesque valores expr. $\surd-1\ (\mathrm{mod.}\ M)$ hi: $N,-N,N',-N',N'',-N''$ etc. Tum per art. 176 forma $(M,N,\frac{NN+1}{M})$ formae $(1,0,1)$ proprie aequivalens erit. Sit transformatio aliqua propria huius in illam, $x=\alpha x'+\beta y'$, $y=\gamma x'+\delta y'$, eruntque repraesentationes numeri $M$ per formam $xx+yy$ ad $N$ pertinentes hae quatuor\HistoricalNote{*)}{Patet enim, hunc casum sub (2) art. {180} contentum esse.}: $x=\pm\alpha,y=\pm\gamma$; $x=\mp\gamma,y=\pm\alpha$.

Quum forma $(1,0,1)$ sit anceps, patet, etiam formam $(M,-N,\frac{NN+1}{M})$ ipsi proprie aequivalentem fore, illamque proprie in hanc transmutari positis $x=\alpha x'-\beta y'$, $y=-\gamma x'+\delta y'$. Hinc derivantur quatuor repraesentationes ipsius $M$ ad $-N$ pertinentes, $x=\pm\alpha,y=\mp\gamma$; $x=\pm\gamma,y=\pm\alpha$. Manifestum itaque est, octo repraesentationes ipsius $M$ dari, quarum semissis altera ad $N$, altera ad $-N$ pertineat; sed hae omnes \textit{unicam} tantummodo discerptionem numeri $M$ in duo quadrata exhibent, $M=\alpha\alpha+\gamma\gamma$, siquidem ad quadrata ipsa tantum, neque vero ad ordinem radicumve signa spectamus.

Quodsi itaque alii valores expr. $\surd-1\ (\mathrm{mod.}\ M)$ praeter $N$ et $-N$ non dantur, quod \textit{e. g.} evenit, quando $M$ est numerus primus, $M$ unico tantum modo in duo quadrata inter se prima resolvi poterit Iam quum $-1$ sit residuum quadraticum cuiusvis numeri primi formae $4n+1$ (art. 108), manifestoque numerus primus in duo quadrata inter se non prima discerpi nequeat, habemus theorema:

\textit{Quivis numerus primus formae $4n+1$ in duo quadrata decomponi potest, et quidem unico tantum modo.}
\[
\begin{gathered}
1=0+1,\quad5=1+4,\quad13=4+9,\quad17=1+16,\\
29=4+25,\quad37=1+36,\quad41=16+25,\quad53=4+49,\\
61=25+36,\quad73=9+64,\quad89=25+64,\quad97=16+81\quad\text{etc.}
\end{gathered}
\]
Theorema hoc elegantissimum iam Fermatio notum fuit. sed ab ill. Eulero primo demonstratum est, \textit{Comm. nov. Petr. T.}\ V, \textit{ad annos} 1754, 1755, \textit{p.}\ 3 sqq. In \textit{T.}\ IV, diss. exstat ad idem argumentum pertinens, \textit{p.}\ 3 sqq., sed tum rem penitus nondum absolverat, vid. imprimis art. 27.

\SourceLeaf{166}
Si igitur numerus aliquis formae $4n+1$ aut pluribus modis aut nullo modo in duo quadrata resolvi potest, certo non erit primus.

Vice versa autem, si expr. $\surd-1\ (\mathrm{mod.}\ M)$ praeter $N$ et $-N$ alios adhuc valores habet, aliae adhuc repraesentationes ipsius $M$ dabuntur, ad hos pertinentes. In hoc itaque casu $M$ pluribus modis in duo quadrata resolvi poterit \textit{e. g.} $65=1+64=16+49$, $221=25+196=100+121$.

Repraesentationes reliquae, in quibus $x,y$ valores obtinent non primos inter se, per methodum nostram generalem facile inveniri possunt. Observamus tantummodo, si numerus aliquis factores formae $4n+3$ involvens, per nullam divisionem per quadratum ab his liberari possit (quod fiet, si aliquis aut plures talium factorum \textit{dimensionem imparem} habet), hunc nullo modo in duo quadrata resolvi posse\HistoricalNote{*)}{Si numerus $M=2^{\mu}Sa^{\alpha}b^{\beta}c^{{\gamma}}\ldots$ ita ut $a,b,c$ etc. sint numeri primi inaequales formae $4n+1$, atque $S$ productum ex omnibus factoribus primis ipsius $M$ formae $4n+3$ (ad quam formam quivis numerus positivus reduci potest, faciendo $\mu=0$ quando $M$ est impar, et $S=1$ quando $M$ nullos factores formae $4n+3$ implicat): $M$ \textit{nullo modo} in duo quadrata resolvi poterit, si $S$ est non-quadratus; si vero $S$ est quadratus, dabuntur $\tfrac12(\alpha+1)(\beta+1)(\gamma+1)$ etc. discerptiones ipsius $M$, quando aliquis numerorum $\alpha,\beta,\gamma$ etc. est impar, aut $\tfrac12(\alpha+1)(\beta+1)(\gamma+1)\text{ etc.}+\tfrac12$, quando omnes $\alpha,\beta,\gamma$ etc. sunt pares (siquidem ad quadrata ipsa tantum respicitur). Qui in calculo combinationum aliquantum sunt versati, demonstrationem huius theorematis (cui, perinde ut aliis \textit{particularibus}. immorari nobis non licet) ex theoria nostra generali haud difficulter eruere poterunt. Cf. art. 105.}.

II. Per formam $xx+2yy$ nullus numerus, cuius non-residuum $-2$. ita repraesentari potest, ut $x$ ad $y$ sit primus, reliqui omnes poterunt. Sit $-2$ residuum numeri $M$, atque $N$ valor aliquis expr. $\surd-2\ (\mathrm{mod.}\ M)$. Tum per art. 176 formae $(1,0,2)$, $(M,N,\frac{NN+2}{M})$ proprie aequivalentes erunt. Transeat illa proprie in hanc ponendo $x=\alpha x'+\beta y'$, $y=\gamma x'+\delta y'$, eritque $x=\alpha,y=\gamma$ repraesentatio numeri $M$ ad $N$ pertinens. Praeter quam et hanc $x=-\alpha,y=-\gamma$ aliae ad $N$ non pertinebunt (art. 180).

Simili modo, ut supra, perspicitur, repraesentationes $x=\pm\alpha,y=\mp\gamma$ ad valorem $-N$ pertinere. Omnes vero hae quatuor repraesentationes unicam tantum discerptionem ipsius $M$ in quadratum et quadratum duplex exhibent, et si praeter $N$ et $-N$ alii valores expr. $\surd-2\ (\mathrm{mod.}\ M)$ non dantur, aliae discerptiones non dabuntur. Hinc adiumento proposs. art. 116 facile deducitur theorema:

\SourceLeaf{167}
\textit{Quivis numerus primus formae $8n+1$ vel $8n+3$ in quadratum et quadratum duplex decomponi potest et quidem unico tantum modo.}
\[
\begin{gathered}
1=1+0,\quad3=1+2,\quad11=9+2,\quad17=9+8,\\
19=1+18,\quad41=9+32.\quad43=25+18,\quad59=9+50,\\
67=49+18,\quad73=1+72,\quad83=81+2,\\
89=81+8,\quad97=25+72\quad\text{etc.}
\end{gathered}
\]
Etiam hoc theorema, uti plura similia, Fermatio innotuit: sed ill. La Grange primus demonstrationem dedit, \textit{Suite des recherches d'Arithmétique}, Nouv. Mém. de l'Ac. de Berlin 1775, p. 323 sqq. Multa ad idem argumentum pertinentia iam ill. Euler absolverat, \textit{Specimen de usu observationum in mathesi pura} Comm. nov. Petr. T. VI p. 185 sqq. Sed demonstratio completa theorematis semper ipsius industriam elusit, p. 220. Conf. etiam diss. in T. VIII (ad annos 1760, 1761), \textit{Supplementum quorundam theorematum arithmeticorum}, sub fin.

III. Per methodum similem demonstratur, quemvis numerum, cuius residuum quadr. sit $-3$, repraesentari posse aut per formam $xx+3yy$, aut per hanc $2xx+2xy+2yy$, ita ut valor ipsius $x$ ad valorem ipsius $y$ sit primus. Quare quum $-3$ sit residuum omnium numerorum primorum formae $3n+1$ (art. 119) manifestoque per formam $2xx+2xy+2yy$ numeri \textit{pares} tantum repraesentari possint: eodem modo ut supra habetur theorema:

\textit{Quivis numerus primus formae $3n+1$ in quadratum et quadratum triplex decomponi potest, et quidem unico tantum modo.}
\[
\begin{gathered}
1=1+0,\quad7=4+3,\quad13=1+12,\quad19=16+3,\\
31=4+27,\quad37=25+12,\quad43=16+27,\quad61=49+12,\\
67=64+3,\quad73=25+48\quad\text{etc.}
\end{gathered}
\]
Demonstrationem huius theorematis ill. Euler primus tradidit in commentatione modo laudata, \textit{Comm. nov. Petr. T.}\ VIII, \textit{p.}\ 105 sqq.

Simili modo ulterius progredi et \textit{e. g.} ostendere possemus, quemvis numerum primum formae $20n+1$, vel $20n+3$, vel $20n+7$, vel $20n+9$ (quippe quorum residuum $-5$) per alterutram formam $xx+5yy$, $2xx+2xy+3yy$ repraesentari posse, et quidem numeros primos formae $20n+1$ et $20n+9$ per priorem, primos formae $20n+3,20n+7$ per posteriorem, nec non dupla primorum formae $20n+1,20n+9$ per formam $2xx+2xy+3yy$, dupla primorum formae $20n+3,20n+7$ per formam $xx+5yy$: sed hanc propositionem infinitasque alias particulares quivis proprio marte ex praecedentibus et infra
\SourceLeaf{168}
tradendis derivare poterit. --- Transimus itaque ad \textit{formas determinantis positivi}, et quum harum indoles prorsus alia sit, quando determinans est quadratus, alia, quando non-quadratus: formas determinantis quadrati hic primo excludimus posteaque seorsim considerabimus.

\Needspace{12\baselineskip}
\begin{center}
\textit{De formis determinantis positivi non-quadrati.}
\end{center}

\Needspace{8\baselineskip}
\Article{183}
\textsc{Problema.} \textit{Proposita forma quacunque $(a,b,a')$, cuius determinans positivus non-quadratus $=D$: invenire formam huic proprie aequivalentem, $(A,B,C)$, in qua $B$ sit positivus et $<\surd D$; $A$ vero si est positivus, vel $-A$, si $A$ negativus, inter $\surd D+B$ et $\surd D-B$ situs.}

\textit{Sol.} Supponimus in forma proposita utramque conditionem nondum locum habere; alioquin enim aliam formam quaerere opus non esset. Porro observamus, in forma determinantis \textit{non-quadrati} terminum primum vel ultimum $=0$ esse non posse (art. 171 ann.). Sit $b'\equiv-b\ (\mathrm{mod.}\ a')$ atque intra limites $\surd D$ et $\surd D\mp a'$ situs (accepto signo superiori quando $a'$ positivus, inferiori quando est negativus) quod fieri posse simili ratione ut art. 3, facile demonstratur, ponaturque $\frac{b'b'-D}{a'}=a''$, qui erit integer, quia $b'b'-D\equiv bb-D\equiv aa'\equiv0\ (\mathrm{mod.}\ a')$. Iam si $a''<a'$, fiat denuo $b''\equiv-b'\ (\mathrm{mod.}\ a'')$ et inter $\surd D$ et $\surd D\mp a''$ situs (prout $a''$ positivus vel negativus) et $\frac{b''b''-D}{{a''}}=a'''$. Si hic iterum $a'''<a''$, sit rursus $b'''\equiv-b''\ (\mathrm{mod.}\ a''')$, et inter $\surd D$ et $\surd D\mp a'''$ situs atque $\frac{b'''b'''-D}{a'''}=a''''$. Haec operatio continuetur, donec in progressione $a',a'',a''',a''''$ etc. ad terminum $a^{m+1}$ perveniatur, praecedente $a^m$ non minorem, quod tandem evenire debet, quia alioquin progressio infinita numerorum integrorum continuo decrescentium haberetur. Tum positis $a^m=A$, $b^m=B$, $a^{m+1}=C$, forma $(A,B,C)$ omnibus conditionibus satisfaciet.

\textit{Dem.} I. Quoniam in progressione formarum $(a,b,a')$, $(a',b',a'')$, $(a'',b'',a''')$ etc. quaevis praecedenti est contigua: ultima $(A,B,C)$ primae $(a,b,a')$ proprie aequivalens erit.

II. Quia $B$ inter $\surd D$ et $\surd D\mp A$ situs est (accipiendo semper signum superius quando $A$ est positivus, inferius quando $A$ est negativus): patet, si ponatur $\surd D-B=p$, $B-(\surd D\mp A)=q$. hos $p,q$ fore positivos. Iam facile confirmatur, fore $qq+2pq+2p\surd D=D+AA-BB$; quare $D+AA-BB$ erit numerus positivus. quem ponemus $=r$. Hinc propter $D=BB-AC$,
\SourceLeaf{169}
fit $r=AA-AC$, adeoque $AA-AC$ numerus positivus: quia vero per hyp. $A$ non est maior quam $C$, manifesto illud aliter fieri nequit, quam si $AC$ est negativus, adeoque signa ipsorum $A,C$ opposita. Hinc $BB=D+AC<D$ adeoque $B<\surd D$.

III. Porro quia $-AC=D-BB$, erit $AC<D$, et hinc (quia $A$ non $>C$), $A<\surd D$. Quare $\surd D\mp A$ erit positivus, adeoque etiam $B$, qui inter limites $\surd D$ et $\surd D\mp A$ est situs.

IV. Hinc a potiori $\surd D+B\mp A$ positivus, et quia $\surd D-B\mp A=-q$. est negativus, $\pm A$ situs erit inter $\surd D+B$ et $\surd D-B$. \textit{Q. E. D.}

\textit{Ex.} Proposita sit forma $(67,97,140)$, cuius determinans $=29$. Hic invenitur progressio formarum $(67,97,140)$, $(140,-97,67)$, $(67,-37,20)$. $(20,-3,-1)$, $(-1,5,4)$. Ultima erit quaesita.

Tales formas $(A,B,C)$ determinantis positivi non-quadrati $D$, in quibus $A$ positive acceptus iacet inter $\surd D+B$ et $\surd D-B$, $B$ vero positivus est atque $<\surd D$, \textit{formas reductas} vocabimus. Formae itaque reductae determinantis positivi non-quadrati aliquantum differunt a formis reductis determinantis negativi; sed propter magnam analogiam inter has et illas, denominationes diversas introducere noluimus.

\Needspace{8\baselineskip}
\Article{184}
Si aequivalentia duarum formarum \textit{reductarum} determinantis positivi aeque facile dignosci posset, ut in formis determinantis negativi (art. 172), aequivalentiam duarum formarum \textit{quarumcunque} eiusdem determinantis positivi nullo negotio diiudicare possemus. Sed hic res longe aliter se habet, fierique potest ut permultae formae reductae inter se aequivalentes sint. Antequam itaque problema hoc aggrediamur, profundius in naturam formarum reductarum (determinantis positivi non-quadrati, quod semper hic subintelligendum) inquirere necesse erit.

1) Si $(a,b,c)$ est forma reducta, $a$ et $c$ signa opposita habebunt. Nam posito determinante formae $=D$, erit $ac=bb-D$, adeoque, propter $b<\surd D$. negativus.

2) Numerus $c$ perinde ut $a$, positive acceptus, inter $\surd D+b$ et $\surd D-b$ situs erit. Nam $-c=\frac{D-bb}{a}$; quare, abstractione facta a signo, $c$ iacebit inter $\frac{D-bb}{\surd D+b}$ et $\frac{D-bb}{\surd D-b}$ \textit{i. e.} inter $\surd D-b$ et $\surd D+b$.

3) Hinc patet, etiam $(c,b,a)$ fore formam reductam.

\SourceLeaf{170}
4) Tum $a$ tum $c$ erunt $<2\surd D$. Uterque enim est $<\surd D+b$, adeoque a potiori $<2\surd D$.

5) Numerus $b$ situs erit inter $\surd D$ et $\surd D\mp a$ (accepto signo superiori quando $a$ positivus, inferiori quando est negativus). Quia enim $\pm a$ iacet inter $\surd D+b$ et $\surd D-b$, erit $\pm a-(\surd D-b)$, sive $b-(\surd D\mp a)$ positivus; $b-\surd D$ autem est negativus; quamobrem $b$ inter $\surd D$ et $\surd D\mp a$ erit situs. --- Prorsus eodem modo demonstratur, $b$ inter $\surd D$ et $\surd D\mp c$ iacere (prout $c$ pos. vel neg.).

6) \textit{Cuivis formae reductae $(a,b,c)$ ab utraque parte contigua est reducta una, et non plures.}

Fiat $a'=c$, $b'\equiv-b\ (\mathrm{mod.}\ a')$ et inter $\surd D$ et $\surd D\mp a'$ situs\HistoricalNote{*)}{Ubi signa ambigua sunt, superiora semper valent quando $a'$ est positivus, inferiora quando $a'$ negativus.}, $c'=\frac{b'b'-D}{a'}$, eritque forma $(a',b',c')$ formae $(a,b,c)$ ab ultima parte contigua, simulque manifestum est, si ulla forma reducta formae $(a,b,c)$ ab ultima parte contigua detur, eam ab hac $(a',b',c')$ diversam esse non posse. Hanc vero revera esse reductam, ita demonstramus.

A) Si ponitur
\[
\surd D+b\mp a'=p,\quad\pm a'-(\surd D-b)=q,\quad\surd D-b=r
\]
hi $p,q,r$ ex (2) supra et defin. formae reductae erunt positivi. Porro ponatur
\[
b'-(\surd D\mp a)=q',\quad\surd D-b'=r'
\]
eruntque $q',r'$ positivi, quia $b'$ iacet inter $\surd D$ et $\surd D\mp a'$. Denique sit $b+b'=\pm ma$ eritque $m$ integer. Iam patet esse $p+q'=b+b'$, adeoque $b+b'$ sive $\pm ma'$ positivum, et proin etiam $m$; unde sequitur $m-1$ certe non esse negativum. Porro fit
\[
r+q'\pm ma'=2b'\pm a',\quad\text{sive}\quad2b'=r+q'\pm(m-1)a'
\]
unde $2b'$ et $b'$ necessario erunt positivi. Et quoniam $b'+r'=\surd D$, erit $b'<\surd D$.

B) Porro fit
\[
r\pm ma'=\surd D+b',\quad\text{sive}\quad r\pm(m-1)a'=\surd D+b'\mp a'
\]
quare $\surd D+b'\mp a'$ erit positivus. Hinc et quoniam $\pm a'-(\surd D-b')=q'$,
\SourceLeaf{171}
adeoque positivus, $\pm a'$ iacebit inter $\surd D+b'$ et $\surd D-b'$. --- Quocirca $(a',b',c')$ erit forma reducta.

Eodem modo demonstratur, si fiat ${}'c=a$, ${}'b\equiv-b\ (\mathrm{mod.}\ {}'c)$ et inter $\surd D$ et $\surd D\pm{}'c$ situs, ${}'a=\frac{{}'b{}'b-D}{{}'c}$, formam $({}'a,{}'b,{}'c)$ fore reductam. Manifesto autem forma haec formae $(a,b,c)$ a parte prima est contigua, aliaque reducta praeter $({}'a,{}'b,{}'c)$ hac proprietate praedita esse non poterit.

\textit{Ex.} Formae reductae $(5,11,-14)$, cuius determinans $=191$, a parte ultima contigua reducta $(-14,3,13)$, a parte prima vero haec $(-22,9,5)$.

7) Si formae reductae $(a,b,c)$ a parte ultima contigua est reducta $(a',b',c')$: reductae $(c,b,a)$ contigua erit a prima parte forma $(c',b',a')$; et si reductae $(a,b,c)$ a prima parte contigua est forma $({}'a,{}'b,{}'c)$; reductae $(c,b,a)$ reducta $({}'c,{}'b,{}'a)$ contigua erit ab ultima parte. Porro etiam formae $(-{}'a,{}'b,-{}'c)$, $(-a,b,-c)$, $(-a',b',-c')$ reductae erunt, et secunda primae, tertia secundae ab ultima parte contiguae, sive prima secundae, secundaque tertiae a parte prima; similiterque tres formae $(-c',b',-a')$, $(-c,b,-a)$, $(-{}'c,{}'b,-{}'a)$. Haec tam obvia sunt ut explicatione non egeant.

\Needspace{8\baselineskip}
\Article{185}
Multitudo omnium formarum reductarum determinantis dati $D$ semper est finita, ipsae vero duplici modo inveniri possunt. Designemus indefinite omnes formas reductas determinantis $D$ per $(a,b,c)$, ita ut omnes valores ipsorum $a,b,c$ determinare oporteat.

\textit{Methodus prima.} Accipiantur pro $a$ omnes numeri (tum positive, tum negative) minores quam $2\surd D$, quorum residuum quadraticum $D$, et pro singulis $a$, ponatur $b$ aequalis omnibus valoribus positivis expr. $\surd D\ (\mathrm{mod.}\ a)$ inter $\surd D$ et $\surd D\mp a$ iacentibus, $c$ vero pro singulis valoribus determinatis ipsorum $a,b$, ponatur $=\frac{bb-D}{a}$. Si quae formae hoc modo oriuntur, in quibus $\pm a$ extra $\surd D+b$ et $\surd D-b$ situs est, reiiciendae sunt.

\textit{Methodus secunda.} Accipiantur pro $b$ omnes numeri positivi minores quam $\surd D$, pro singulis $b$ resolvatur $bb-D$ omnibus quibus fieri potest modis in binos factores, qui neglecto signo inter $\surd D+b$ et $\surd D-b$ iaceant, ponaturque alter $=a$, alter $=c$. Manifestum est, singulas resolutiones in factores praebere binas formas, quia uterque factor tum $=a$, tum $=c$ poni debet.
